
\documentclass[11pt]{article}
\usepackage{series}
\usepackage{array}

\hypersetup{
  pdftitle={Gauge Fixing as Admissible Section Selection in Modal Triplet Theory},
  pdfauthor={Peter Nero},
  pdfsubject={Modal Triplet Theory, gauge fixing, admissible projection, Faddeev--Popov, BRST},
  pdfkeywords={Modal Triplet Theory, gauge fixing, projection, Faddeev-Popov determinant, BRST, Gribov ambiguity, Dirac delta}
}

\title{\textbf{Gauge Fixing as Admissible Section Selection in Modal Triplet Theory}\\
\large Faddeev--Popov Determinants, Ghosts, BRST, and Gribov Ambiguity as Projection Geometry}
\author{Peter Nero}
\date{July 2026, Version 1}

\newcommand{\A}{\mathcal A}
\newcommand{\G}{\mathcal G}
\newcommand{\M}{\mathcal M}
\newcommand{\gh}{\mathfrak g}
\newcommand{\FP}{\mathrm{FP}}
\newcommand{\adm}{\mathrm{adm}}
\newcommand{\coh}{\mathrm{coh}}
\newcommand{\vol}{\mathrm{vol}}
\newcommand{\id}{\mathrm{id}}
\newcommand{\Lie}{\mathcal L}

\begin{document}
\maketitle

\begin{abstract}
Gauge theory contains an explicit version of the projection architecture developed in Modal Triplet Theory (MTT). The full field space \(\A\) contains redundant representatives related by a gauge group \(\G\), while the physical configuration space is the quotient \(\A/\G\). A gauge condition \(G[A]=0\) is therefore not a physical law imposed on the quotient itself, but a local representative section of the quotient map. This paper develops the MTT interpretation of gauge fixing as admissible section selection. The gauge-fixing delta \(\delta(G[A])\) is identified as the singular zero-width limit of a representative-selection kernel, the Faddeev--Popov determinant is the Jacobian of the projection from gauge-orbit coordinates to slice coordinates, ghost fields are quotient-measure bookkeeping variables, BRST cohomology is the algebraic implementation of admissible quotienting, and Gribov ambiguity is the failure of a global admissible section.

The mathematical core is deliberately modest. We first prove a finite-dimensional slice formula showing that the Faddeev--Popov determinant is exactly the transverse projection Jacobian for a local quotient. We then state the corresponding regularized infinite-dimensional functional identity under the usual formal assumptions of gauge-fixed path integration. Finally, we give a typed MTT interpretation of these standard structures. The section language is exact, but the assignment of lens, circle, and nil roles is an encoding dictionary rather than a derivation of gauge theory. A gauge section chooses a representative of an orbit; it neither inverts the quotient globally nor recovers a unique ontic configuration. The actual global obstruction depends on the topology and orbit stratification of the declared gauge bundle.
\end{abstract}

\section*{Version 1 Revision Note}
\addcontentsline{toc}{section}{Version 1 Revision Note}
\begin{description}
\item[Supersedes] The unversioned April 2026 manuscript with the same title.
\item[Reason] The local-section language was sound, but the text did not
state sharply enough that a right inverse chooses a representative rather
than recovering unique ontology, and it treated the MTT triadic dictionary
as more than an interpretation of standard gauge geometry.
\item[Resolution] This version separates the rigorous finite-dimensional
slice theorem, the regulated formal field-theory identity, and the MTT
encoding.  It records free-action and transversality hypotheses, orbit
strata and stabilizers, and the genuine global-section obstruction.
\item[Retained result] Locally, a transverse gauge condition defines a
representative section and its Faddeev--Popov determinant is the associated
orbit-to-slice Jacobian.
\item[Remaining boundary] The paper does not construct a global gauge
section, a nonperturbative gauge measure, or an MTT-selected BV/BRST
quantization.
\end{description}

\seriestagline

\tableofcontents

\section{Purpose and claim discipline}

The previous paper in this sequence argued that Dirac delta functions should be read, in MTT-compatible effective descriptions, as singular shadows of finite admissible projection kernels. The present paper applies that principle to the most transparent standard case: gauge fixing.

The guiding slogan is:
\[
\boxed{\text{Gauge is uncollapsed projection; gauge fixing is collapsed representative selection.}}
\]

More explicitly, gauge freedom records that the map from a richer description to a physical quotient is many-to-one. Gauge fixing chooses a representative section. The Faddeev--Popov determinant corrects the measure induced by that choice. The gauge-fixing delta is the zero-width limit of a finite representative-selection kernel.

\subsection{What is proved}

This paper proves a finite-dimensional local slice formula. In a finite-dimensional free gauge action with a transverse gauge condition, the Faddeev--Popov determinant is exactly the Jacobian by which orbit coordinates project onto the gauge-fixing condition. This is a rigorous local theorem.

\subsection{What is not proved}

This paper does not provide a new rigorous construction of nonperturbative Yang--Mills theory. It does not solve the Gribov problem. It does not replace the BV--BRST formalism. It does not claim that every gauge theory has a global smooth gauge section. In fact, one of the central claims is that global section failure is structurally expected.

The infinite-dimensional functional-integral statements are used in the standard physicists' sense: as formal identities that become rigorous only after choosing a regulator, a measure construction, and a renormalization scheme. The finite-dimensional theorem is the analytic anchor; the functional gauge-theory interpretation is the physical extension.

\begin{center}
\fbox{\begin{minipage}{0.92\textwidth}
\textbf{Scope of proof.} This paper proves the finite-dimensional local slice/Jacobian identity. It then uses the standard Faddeev--Popov formalism as the field-theoretic extension of that identity. The MTT contribution is the interpretation: gauge fixing is admissible representative selection, the gauge-fixing delta is the zero-width selection kernel, and the Faddeev--Popov determinant is the projection Jacobian.
\end{minipage}}
\end{center}

\section{Projection structure of gauge theory}

Let \(\A\) be a field configuration space and let \(\G\) be a gauge group acting on it:
\[
\A\times\G\to\A,\qquad (A,g)\mapsto A^g.
\]
The physical configuration associated with \(A\) is not \(A\) itself but its orbit:
\[
[A]=\{A^g:g\in\G\}.
\]
The physical configuration space is the quotient
\[
\M:=\A/\G.
\]
Thus there is a quotient projection
\[
\pi:\A\to\M,\qquad A\mapsto [A].
\]

\begin{definition}[Gauge orbit]
The gauge orbit through \(A\in\A\) is
\[
\mathcal O_A:=\pi^{-1}([A])=\{A^g:g\in\G\}.
\]
\end{definition}

\begin{definition}[Local gauge section]
A local gauge section over \(U\subset\M\) is a map
\[
s:U\to\A
\]
such that
\[
\pi\circ s=\id_U.
\]
Equivalently, \(s\) chooses one representative \(A=s([A])\) from each orbit in \(U\).
\end{definition}

\begin{remark}[A section is not ontic recovery]
The identity \(\pi\circ s=\id_U\) is a right-inverse relation.  It says that
the selected representative projects back to the same gauge orbit.  It does
not make \(\pi\) injective, reconstruct all representatives, or identify
\(s([A])\) as a unique ontic field configuration.  Different admissible
sections can encode the same physical orbit.
\end{remark}

\begin{remark}[MTT reading]
The quotient projection \(\pi:\A\to\A/\G\) is a lens-type projection. Its fibers are equivalence classes of redundant representatives. Gauge freedom is not additional physical content; it is the visible non-injectivity of the representation map.
\end{remark}

The central structural point is:
\[
\boxed{\text{Gauge freedom}=\text{visible persistence of projection non-injectivity}.}
\]

\section{Gauge conditions as local sections}

A gauge condition is a functional
\[
G:\A\to V
\]
into some vector space \(V\), with the intended gauge slice
\[
\Sigma_G:=\{A\in\A:G[A]=0\}.
\]

If \(\Sigma_G\) intersects each nearby gauge orbit exactly once and transversely, then it defines a local section of the quotient map.

\begin{definition}[Admissible local gauge condition]
A gauge condition \(G[A]=0\) is locally admissible at \(A\) if:
\begin{enumerate}
\item the orbit through \(A\) intersects \(\Sigma_G\);
\item the intersection is locally unique;
\item the intersection is transverse, i.e. the linearized map from infinitesimal gauge directions to \(V\) is invertible.
\end{enumerate}
\end{definition}

The linearized gauge variation is
\[
\delta_\epsilon A := R_A\epsilon,
\]
where \(R_A:\gh\to T_A\A\) is the infinitesimal gauge-action map. The gauge-fixing operator is
\[
M_G(A):=D_A G\circ R_A:\gh\to V.
\]
The transverse condition is:
\[
\det M_G(A)\ne0.
\]

\begin{remark}
In ordinary gauge theory \(M_G(A)\) is the Faddeev--Popov operator. For example, in Yang--Mills theory with Lorenz gauge \(G[A]=\partial^\mu A_\mu\), the infinitesimal variation gives an operator of the schematic form
\[
M_G(A)=\partial^\mu D_\mu[A].
\]
\end{remark}

\section{Finite-dimensional slice theorem}

We now state the rigorous finite-dimensional model behind the Faddeev--Popov identity.

\begin{assumption}[Finite-dimensional local quotient model]
Let \(X\) be a smooth oriented \(n\)-dimensional manifold, and let a \(k\)-dimensional Lie group \(G\) act freely and properly on \(X\). Let \(G_0:X\to\mathbb R^k\) be a smooth gauge condition. Fix \(x_0\in X\) such that:
\begin{enumerate}
\item \(G_0(x_0)=0\);
\item the differential \(D_{x_0}G_0\) restricted to the tangent space of the orbit \(T_{x_0}\mathcal O_{x_0}\) is an isomorphism onto \(\mathbb R^k\).
\end{enumerate}
\end{assumption}

Let \(R_{x_0}:\mathfrak g\to T_{x_0}X\) be the infinitesimal action map. Define
\[
M_{G_0}(x_0):=D_{x_0}G_0\circ R_{x_0}:\mathfrak g\to\mathbb R^k.
\]

\begin{theorem}[Local slice and projection Jacobian]
Under the finite-dimensional local quotient assumptions, there is a neighborhood \(U\) of \(x_0\) such that \(G_0^{-1}(0)\cap U\) is a local slice through the \(G\)-orbits. Moreover, in local coordinates \((u,\alpha)\), where \(u\) parameterizes the slice and \(\alpha\in\mathfrak g\) parameterizes the orbit, the volume form factorizes as
\[
d\vol_X = J(u,\alpha)\,du\,d\alpha,
\]
and the distributional identity
\[
1
=
\left|\det M_{G_0}(x)\right|
\int_{\mathfrak g} d\alpha\,\delta(G_0(x^\alpha))
\]
holds locally, up to the conventional normalization of Haar measure.
\end{theorem}

\begin{proof}
By the transversality assumption, \(D_{x_0}G_0\) restricted to the orbit tangent space is an isomorphism. The implicit function theorem implies that \(G_0^{-1}(0)\) is a codimension-\(k\) submanifold through \(x_0\), transverse to the orbit. Since the action is free and proper, the orbit map is locally an embedding, and the map
\[
(u,\alpha)\mapsto u^\alpha
\]
from slice coordinates and group coordinates is a local diffeomorphism onto a neighborhood of \(x_0\).

For fixed \(u\), define
\[
F_u(\alpha):=G_0(u^\alpha).
\]
At the selected representative \(\alpha=0\), the derivative is
\[
D_\alpha F_u(0)=D_uG_0\circ R_u=M_{G_0}(u).
\]
By invertibility, the ordinary multidimensional delta-change-of-variables formula gives
\[
\delta(F_u(\alpha))
=
\frac{\delta(\alpha)}{|\det M_{G_0}(u)|}
\]
locally. Therefore
\[
\int d\alpha\,\delta(G_0(u^\alpha))
=
\frac{1}{|\det M_{G_0}(u)|}.
\]
Multiplying by \(|\det M_{G_0}(u)|\) gives the stated identity. Extending from \(u\) to \(x\) in the local neighborhood gives the same expression for \(x\) near the slice. \end{proof}

\begin{corollary}[Faddeev--Popov determinant as projection Jacobian]
In a local gauge slice, the Faddeev--Popov determinant is the Jacobian measuring how gauge-orbit coordinates project onto the gauge-fixing condition:
\[
\Delta_{\FP}(x)=|\det(DG_0\circ R_x)|.
\]
\end{corollary}

\begin{remark}[MTT interpretation]
The determinant is not mysterious. It is the measure cost of reducing a redundant fiber to a chosen local representative. In MTT language, it is the Jacobian of lens-quotient projection.
\end{remark}

\section{Soft representative selection}

The delta function in the Faddeev--Popov identity imposes a perfectly sharp gauge slice:
\[
\delta(G[A]).
\]
The delta-projection principle replaces such a sharp condition with a finite admissibility kernel. For a scale \(\epsilon>0\), define
\[
\mathcal K_\epsilon(G[A])
:=
(2\pi\epsilon^2)^{-k/2}
\exp\left(-\frac{\|G[A]\|^2}{2\epsilon^2}\right)
\]
in finite dimension, or the corresponding Gaussian functional weight in field theory.

Then
\[
\mathcal K_\epsilon(G[A])\to \delta(G[A])
\]
distributionally as \(\epsilon\downarrow0\).

\begin{definition}[Gauge tube]
For \(\epsilon>0\), the \(\epsilon\)-gauge tube around the slice \(G[A]=0\) is the finite-width representative-selection region weighted by \(\mathcal K_\epsilon(G[A])\).
\end{definition}

\begin{proposition}[Soft slice identity]
In the finite-dimensional local quotient model,
\[
\Delta_{\FP}(x)\int d\alpha\,\mathcal K_\epsilon(G(x^\alpha))
\to 1
\]
as \(\epsilon\downarrow0\), locally near a transverse slice.
\end{proposition}

\begin{proof}
This follows directly from the distributional convergence \(\mathcal K_\epsilon\to\delta\) and the local slice theorem. \end{proof}

\begin{remark}[MTT reading]
A hard gauge slice is an idealized zero-width representative choice. A finite gauge tube is the MTT-native version: representative selection with finite admissibility width.
\end{remark}

\section{Functional Faddeev--Popov identity}

In field theory one writes the formal identity
\[
1=\Delta_{\FP}[A]\int\mathcal D\alpha\,\delta(G[A^\alpha]),
\]
where
\[
\Delta_{\FP}[A]=\det\left(\frac{\delta G[A^\alpha]}{\delta\alpha}\right)_{\alpha=0}.
\]

In infinite dimension this determinant is formal until a regulator, boundary condition, and renormalization prescription are specified. In what follows it is used in the standard perturbative sense. The rigorous content imported from the previous section is the finite-dimensional local slice formula; the field-theoretic expression is its regulated formal analogue.

Gauge-fixed path integration then takes the form
\[
Z
=
\int \mathcal DA\,\Delta_{\FP}[A]\,\delta(G[A])\,e^{iS[A]/\hbar}.
\]

The MTT-consistent soft form is
\[
Z_{\adm,\epsilon}
=
\int \mathcal DA\,\Delta_{\FP}[A]\,
\mathcal K_\epsilon(G[A])\,
e^{iS[A]/\hbar}.
\]

For a quadratic Gaussian gauge-tube kernel,
\[
\mathcal K_\epsilon(G[A])
=
\exp\left(-\frac{1}{2\epsilon^2}\|G[A]\|^2\right),
\]
the familiar covariant gauge-fixing term appears as the finite-width representative-selection penalty.

\begin{remark}
This is not a new gauge theory. It is a reinterpretation of the standard gauge-fixing term: the gauge parameter controls the width of the representative-selection tube. The exact delta gauge is recovered as the singular zero-width limit.
\end{remark}

\section{Worked examples: Abelian and non-Abelian Lorenz gauge}

The abstract projection-Jacobian statement becomes especially transparent in the standard Lorenz gauge.

\subsection{\texorpdfstring{Abelian \(U(1)\)}{Abelian U(1)}}

Let \(A_\mu\) be an Abelian gauge potential and let
\[
A_\mu^\alpha=A_\mu+\partial_\mu\alpha .
\]
Choose the Lorenz gauge condition
\[
G[A]=\partial^\mu A_\mu .
\]
Then
\[
G[A^\alpha]
=
\partial^\mu A_\mu+\Box\alpha .
\]
Therefore the Faddeev--Popov operator is
\[
M_G=\Box .
\]
Hence
\[
\Delta_{\FP}[A]=\det \Box ,
\]
which is independent of \(A\). In the Abelian case the ghost determinant is therefore field-independent and decouples from the interacting dynamics.

\begin{remark}[MTT reading]
For \(U(1)\), the local projection Jacobian is constant along the physical field directions. The quotient still exists, and the gauge-fixing delta still selects a representative slice, but the measure cost does not introduce interacting quotient-bookkeeping fields.
\end{remark}

The soft MTT gauge tube gives the familiar covariant gauge penalty
\[
\exp\left[-\frac{1}{2\epsilon^2}\int(\partial^\mu A_\mu)^2\right],
\]
in Euclidean signature, or the corresponding phase-weighted term in Lorentzian signature. Thus the usual gauge parameter can be read as controlling the width of admissible representative selection around the Lorenz slice.

\subsection{Non-Abelian Yang--Mills}

For a non-Abelian gauge field \(A_\mu=A_\mu^aT_a\), the infinitesimal gauge variation is
\[
\delta_\alpha A_\mu = D_\mu[A]\alpha ,
\]
where
\[
D_\mu[A]\alpha=\partial_\mu\alpha+[A_\mu,\alpha].
\]
With the same Lorenz gauge condition,
\[
G[A]=\partial^\mu A_\mu,
\]
one obtains
\[
G[A^\alpha]=G[A]+\partial^\mu D_\mu[A]\alpha+O(\alpha^2).
\]
Therefore
\[
M_G[A]=\partial^\mu D_\mu[A].
\]
The Faddeev--Popov determinant is now
\[
\Delta_{\FP}[A]=\det(\partial^\mu D_\mu[A]),
\]
which depends on the gauge field. Consequently the ghost representation produces interacting ghost terms.

\begin{remark}[Why ghosts interact]
Ghosts interact in the non-Abelian case because the projection Jacobian depends on the point of the gauge orbit. Equivalently, the redundancy fiber is curved or twisted relative to the chosen representative slice. Ghost interactions therefore encode the field-dependent geometry of quotienting, not additional physical matter.
\end{remark}

\begin{center}
\begin{tabular}{
>{\raggedright\arraybackslash}p{0.26\textwidth}
>{\raggedright\arraybackslash}p{0.28\textwidth}
>{\raggedright\arraybackslash}p{0.36\textwidth}}
\hline
Case & Faddeev--Popov operator & Interpretation\\
\hline
Abelian \(U(1)\) & \(\Box\) & constant projection Jacobian; ghosts decouple\\
Non-Abelian Yang--Mills & \(\partial^\mu D_\mu[A]\) & field-dependent projection Jacobian; ghosts interact\\
\hline
\end{tabular}
\end{center}

\section{Ghosts as quotient bookkeeping}

For non-Abelian gauge theory the Faddeev--Popov determinant is field-dependent. One introduces anticommuting ghost fields \(c,\bar c\) to represent it:
\[
\Delta_{\FP}[A]
=
\int \mathcal D\bar c\,\mathcal Dc\,
\exp\left(iS_{\mathrm{ghost}}[\bar c,c,A]\right).
\]

In the present interpretation:
\[
\boxed{\text{ghost fields are local bookkeeping variables for the quotient-measure Jacobian.}}
\]

They are not physical excitations in the same sense as gauge-invariant particles. They encode how the redundant lens directions distort the measure after representative selection.

\begin{proposition}[Ghosts encode the determinant]
Formally, for a linearized Faddeev--Popov operator \(M_G[A]\),
\[
\det M_G[A]
=
\int \mathcal D\bar c\,\mathcal Dc\,
\exp\left(-\int \bar c\,M_G[A]\,c\right)
\]
in Euclidean signature, up to normalization.
\end{proposition}

\begin{remark}
The sign and factor of \(i\) depend on signature and convention. The structural point is invariant: ghost integration represents the determinant induced by quotienting the gauge orbit.
\end{remark}

\section{BRST as algebraic admissible quotienting}

BRST symmetry packages gauge redundancy, gauge fixing, and ghosts into a cohomological structure. The BRST operator \(Q_{\mathrm{BRST}}\) satisfies
\[
Q_{\mathrm{BRST}}^2=0.
\]
Physical states are represented by cohomology:
\[
\mathcal H_{\mathrm{phys}}
=
\ker Q_{\mathrm{BRST}}/\operatorname{im}Q_{\mathrm{BRST}}.
\]

In MTT language:
\[
\boxed{
\text{physical states}=
\text{admissible quotient classes modulo null redundancy}.
}
\]

Exact BRST states are invisible redundancy; closed BRST states are compatible with the quotient constraint.

\begin{remark}[Nilpotency and projection]
The nilpotency \(Q_{\mathrm{BRST}}^2=0\) says that gauge redundancy has already been identified as a null direction of physical distinction. Acting twice remains within redundancy. This is the algebraic counterpart of a projection quotient.
\end{remark}

\section{Gribov ambiguity as global section failure}

The local slice theorem does not imply global gauge fixing. A gauge condition may intersect a single orbit more than once:
\[
A_1,A_2\in\Sigma_G,\qquad A_2=A_1^g,\qquad A_1\ne A_2.
\]
These are Gribov copies.

In the present language:
\[
\boxed{
\text{Gribov ambiguity}=
\text{failure of a global admissible representative section}.
}
\]

This is not an accidental technical nuisance. It is a structural warning:
the quotient projection may have nontrivial global topology, and no single
downstream chart can cover it without ambiguity.  For a free action this is
the global-section problem for the principal gauge bundle
\(\A\to\A/\G\).  When stabilizers occur, the quotient is stratified and the
principal-bundle picture applies only orbit-type by orbit-type.  The
Faddeev--Popov operator can also lose invertibility at a Gribov horizon, so a
local slice theorem cannot be promoted to a global gauge choice.

\begin{remark}[MTT alignment]
MTT expects reduced descriptions to be local encoding regimes with controlled overlap, not global primitive ontologies. The Gribov problem is therefore a standard gauge-theoretic example of a general MTT principle: representative choice can be locally admissible while globally obstructed.
\end{remark}

\section{Triadic placement as an encoding dictionary}

The MTT proto-spinorial carrier is written schematically as
\[
\Xi=(\Psi,C,L,N),
\]
where \(C\) carries return/bookkeeping, \(L\) carries lens redundancy transport, and \(N\) carries nil termination or survivor selection.

The following table is a proposed typed encoding of standard gauge-theory
objects.  It does not derive the gauge group, BRST differential, or measure
from the proto-spinor:

\begin{center}
\begin{tabular}{p{0.28\textwidth}p{0.25\textwidth}p{0.37\textwidth}}
\hline
Gauge-theory object & MTT role & Interpretation\\
\hline
Gauge orbit & \(L\) & lens redundancy fiber\\
Gauge fixing & \(L\)-selection interface & representative selection\\
\(\delta(G[A])\) & selection idealization & sharp slice selection\\
Faddeev--Popov determinant & lens Jacobian & quotient-measure cost\\
Ghosts & lens bookkeeping & determinant variables\\
BRST cohomology & quotient algebra & admissible physical classes\\
Gribov copies & section obstruction & local charts only\\
Ward identities & \(C/L\) bookkeeping & return/consistency constraints\\
\hline
\end{tabular}
\end{center}

The key point is that gauge freedom is primarily lens structure, while gauge fixing introduces a representative-selection operation with the same formal profile as nil-sector selection. The delta function appears at precisely this interface:
\[
\boxed{
\text{gauge delta}=\text{singular representative-selection kernel on a lens quotient}.
}
\]

\section{Diagnostic consequences}

The present paper gives a practical diagnostic:

\begin{enumerate}
\item Find the gauge redundancy.
\item Identify the quotient projection.
\item Identify the local gauge section.
\item Compute the projection Jacobian.
\item If a finite representative-selection kernel is introduced, treat it
as a declared gauge-fixing regularization and verify that the BRST/BV
identities and gauge-parameter independence survive.
\item Check whether the local section extends globally; if not, Gribov-type obstruction is expected.
\end{enumerate}

This suggests several later developments:

\begin{enumerate}
\item soft gauge fixing as finite admissible section selection;
\item BRST/BV as the algebra of admissible quotienting;
\item Gribov regions as admissible chart domains;
\item gauge-parameter dependence as tube-width dependence;
\item anomalies as failures of quotient-measure consistency.
\end{enumerate}

\section{Conclusion}

Gauge theory provides a clean standard example of quotient and section
geometry.  The full field space contains redundant representatives; the
physical space is a quotient; gauge fixing chooses a local representative;
the Faddeev--Popov determinant is the local orbit-to-slice Jacobian; ghosts
represent that determinant in the perturbative functional formalism; BRST
cohomology encodes gauge-invariant classes; and Gribov ambiguity obstructs a
single global section.

Thus the central conclusion is:
\[
\boxed{
\begin{minipage}{0.86\textwidth}
\centering
Gauge fixing is admissible section selection, and Faddeev--Popov is its projection Jacobian.
\end{minipage}
}
\]

The MTT contribution is interpretive: its projection and section vocabulary
can organize these established structures without replacing them.  A finite
gauge-selection kernel or an MTT-sourced BRST complex remains a separate
construction, not a consequence of the local slice theorem.

\begin{thebibliography}{9}
\bibitem{FaddeevPopov1967}
L.~D. Faddeev and V.~N. Popov,
\emph{Feynman diagrams for the Yang--Mills field},
Physics Letters B \textbf{25} (1967), 29--30.

\bibitem{Gribov1978}
V.~N. Gribov,
\emph{Quantization of non-Abelian gauge theories},
Nuclear Physics B \textbf{139} (1978), 1--19.

\bibitem{Singer1978}
I.~M. Singer,
\emph{Some remarks on the Gribov ambiguity},
Communications in Mathematical Physics \textbf{60} (1978), 7--12.

\bibitem{HenneauxTeitelboim1992}
M. Henneaux and C. Teitelboim,
\emph{Quantization of Gauge Systems},
Princeton University Press, 1992.
\end{thebibliography}
\end{document}
